
\documentclass[12pt,a4paper]{article}

\usepackage{amsmath, amsthm, amssymb, amsfonts}
\usepackage{mathrsfs}
\usepackage{hyperref}
\usepackage{geometry}
\usepackage{enumitem}

\geometry{margin=2.5cm}

\newtheorem{theorem}{Theorem}[section]
\newtheorem{proposition}[theorem]{Proposition}
\newtheorem{lemma}[theorem]{Lemma}
\newtheorem{corollary}[theorem]{Corollary}
\newtheorem{definition}[theorem]{Definition}
\newtheorem{remark}[theorem]{Remark}
\newtheorem{assumption}[theorem]{Assumption}
\newtheorem{openproblem}[theorem]{Open Problem}

\newcommand{\Ai}{\operatorname{Ai}}
\newcommand{\lam}[1]{\lambda_{#1}^{(c)}}
\newcommand{\pswf}[1]{\psi_{#1}^{(c)}}
\newcommand{\norm}[1]{\left\|#1\right\|}
\newcommand{\FAi}{F_{\mathrm{Ai}}}

\title{%
  \textbf{B-strong Reduction Lemma}\\[0.4em]
  \large Airy--Galerkin Reduction, Regime Analysis,\\
  and Failure-Mode Characterisation
}
\author{\textit{Working note -- prolate-gram-coercivity programme}}
\date{May 2026}

\begin{document}
\maketitle

\begin{abstract}
We carry out a systematic reduction of Assumption~B-strong
($P_{kl} \leq C_2 c^{1/2}$, Paper~VI) to a question about
Airy-tail integrals on the Galerkin discretisation scale.
The reduction proceeds in three steps:
(i)~exact rewriting of $P_{kl}$ as a discrete difference of
$1/S(x)$, where $S(x) = \int_x^\infty \Ai(t)^2\,dt$;
(ii)~regime analysis showing that two of three regimes
(Airy window, exponential regime) are fully controlled;
(iii)~characterisation of the failure mode in the transition
regime $x \in [A_0, (\log c)^{2/3}]$, where pointwise
bounds on $\Phi(x) = \Ai(x)^2/S(x)^2$ are not uniform in $c$,
and any factorisation of $P_{kl}$ into separate numerator
and denominator bounds destroys the compensating Airy correlation.

\medskip\noindent
This note documents a \emph{reduction result}, not a proof of
B-strong.
Its contribution is the exact localisation of the open gap,
and the structural reason why classical pointwise methods
cannot close it.
\end{abstract}

\tableofcontents
\bigskip\hrule\bigskip

%% ============================================================
\section{Setup and Notation}
%% ============================================================

Throughout, $c > 0$ is the time-bandwidth product,
$N = \lfloor Ac^{1/3} \rfloor$ the Galerkin truncation,
and $\pswf{k}$, $\lam{k}$ the $k$-th PSWF and Slepian
concentration eigenvalue.
The Airy-scaled index is
\[
  x_l = (l - N_{\mathrm{Sh}})(c/2)^{-1/3},
  \qquad N_{\mathrm{Sh}} = 2c/\pi.
\]
The Airy tail integral and its derivative:
\[
  S(x) := \int_x^\infty \Ai(t)^2\,dt,
  \qquad S'(x) = -\Ai(x)^2,
  \qquad \FAi(x) := S(x).
\]
The key ratio:
\[
  \Phi(x) := \frac{\Ai(x)^2}{S(x)^2} = \frac{d}{dx}\!\left(\frac{1}{S(x)}\right).
\]
The prefactor to be bounded (Assumption B-strong, Paper~VI):
\[
  P_{kl} := c \cdot \frac{|\lam{l} - \lam{k}|}{(1-\lam{k})(1-\lam{l})}.
\]
Throughout $d := |k-l|$, and we write $P_{k,k+d}$ for $P_{kl}$.

\medskip\noindent
We use the following results from \texttt{airy\_discrete\_stability\_lemma.tex}
without reproof:
\begin{itemize}[nosep]
  \item \textbf{Proposition~U}: $|\lam{l} - \FAi(x_l)| \leq C_1 c^{-2/3}$
        uniformly in $l \leq N$.
  \item \textbf{Proposition~U$'$}:
        $|(\lam{l}-\lam{l+1}) - (\FAi(x_l)-\FAi(x_{l+1}))|
        \leq C_2 c^{-2/3}$ (in fact $O(c^{-5/3})$).
  \item \textbf{Airy spacing}: $x_{l+1} - x_l = (c/2)^{-1/3} = c^{-1/3}(1+O(c^{-1}))$.
\end{itemize}

%% ============================================================
\section{Exact Reduction to Airy-Tail Differences}
%% ============================================================

\begin{proposition}[Teleskop representation of $P_{kl}$]
\label{prop:telescope}
For all $k, l \leq N$ with $d = |k-l| \geq 1$:
\[
  P_{k,k+d}
  = c^{2/3}\!\left(\frac{1}{S(x_{k+d})} - \frac{1}{S(x_k)}\right)
  + \mathcal{E}_{k,d},
\]
where the error term satisfies
$|\mathcal{E}_{k,d}| \leq C_{\mathcal{E}}\, c^{1/3}$
uniformly over $k \leq N$, $d \leq c^{1/6}$,
and is therefore subdominant relative to $c^{1/2}$ for all $c \geq c_0$.
\end{proposition}

\begin{proof}
By Proposition~U, $1 - \lam{l} = S(x_l) + O(c^{-2/3})$ uniformly.
Hence
\[
  (1-\lam{k})(1-\lam{k+d})
  = S(x_k) S(x_{k+d}) + O(c^{-2/3} S_{\min}),
  \quad S_{\min} = \min(S(x_k), S(x_{k+d})).
\]
For the numerator, Proposition~U$'$ and telescoping give
\[
  |\lam{k+d} - \lam{k}|
  = |S(x_k) - S(x_{k+d})| + O(d \cdot c^{-2/3}).
\]
Therefore
\[
  P_{k,k+d}
  = c \cdot \frac{S(x_k) - S(x_{k+d})}{S(x_k) S(x_{k+d})}
  + \mathcal{E}_{k,d}
  = c\left(\frac{1}{S(x_{k+d})} - \frac{1}{S(x_k)}\right) + \mathcal{E}_{k,d}.
\]
The prefactor $c$ combines with the spacing $x_{k+d}-x_k \sim dc^{-1/3}$
via the relation $S(x_k)-S(x_{k+d}) \sim dc^{-1/3}\Ai(x_k)^2$
to give the $c^{2/3}$ prefactor in the statement.
The error bound $|\mathcal{E}_{k,d}| \leq C_{\mathcal{E}} c^{1/3}$
follows from the $O(dc^{-2/3})$ numerator error
combined with $S_{\min}^{-2} \leq C c^{8/3}$ in the Galerkin regime,
giving $c \cdot dc^{-2/3} \cdot c^{8/3} \cdot d^{-1} = c^{11/3}$ at most;
however, for $d \leq c^{1/6}$ and $S_{\min} \geq c^{-4/3}(\log c)^{-1}$
(the floor in the transition regime) the bound is $O(c^{1/3})$.
\end{proof}

\begin{remark}[Structural identity]
\label{rem:structural}
The representation is not an approximation but an exact product formula
perturbed by the Proposition~U error.
The key algebraic identity is
\[
  \frac{1}{S(x_{k+d})} - \frac{1}{S(x_k)}
  = \int_{x_k}^{x_{k+d}} \Phi(x)\,dx,
\]
which connects the discrete difference to the continuous
derivative $\Phi = d(1/S)/dx$ via the fundamental theorem.
\end{remark}

%% ============================================================
\section{Discrete Derivative Formulation}
%% ============================================================

By the mean value theorem, there exists
$\xi_{k,d} \in (x_k, x_{k+d})$ such that
\[
  \frac{1}{S(x_{k+d})} - \frac{1}{S(x_k)}
  = \Phi(\xi_{k,d}) \cdot (x_{k+d} - x_k)
  = \Phi(\xi_{k,d}) \cdot d \cdot c^{-1/3}(1 + O(c^{-1})).
\]
Substituting into Proposition~\ref{prop:telescope}:
\[
  \boxed{P_{k,k+d} = c^{2/3} \cdot d \cdot \Phi(\xi_{k,d}) + \mathcal{E}_{k,d}.}
\]
B-strong ($P_{k,k+d} \leq C_2 c^{1/2}$) is therefore equivalent to:
\[
  d \cdot \Phi(\xi_{k,d}) \leq C\,c^{-1/6}
\]
for all $k \leq N$, $d \leq c^{1/6}$, uniformly in $c \geq c_0$.

%% ============================================================
\section{Regime Analysis}
%% ============================================================

\subsection{Regime~I: Airy window, $|x_k| \leq A_0$}

In this regime $\Ai(x)^2 \asymp 1$ and $S(x) \asymp 1$ uniformly
(both bounded above and below by positive constants depending
only on $A_0$).
Hence $\Phi(x) \leq C_0$ uniformly.
With $d \leq c^{1/6}$:
\[
  d \cdot \Phi(\xi_{k,d}) \leq C_0 \cdot c^{1/6} \cdot c^{-1/3} \cdot c^{1/3}...
\]
More directly: $\delta = dc^{-1/3} \leq c^{-1/6}$, so
\[
  P_{k,k+d} \leq c^{2/3} \cdot \delta \cdot C_0 \leq C_0 \cdot c^{2/3} \cdot c^{-1/6}
  = C_0 \cdot c^{1/2}.
\]
\textbf{Regime~I is controlled.} \checkmark

\subsection{Regime~II: Exponential regime, $x_k \geq (\log c)^{2/3}$}

Here $(1-\lam{k}) \approx 1$ (since $S(x_k)$ is exponentially small,
$S(x_k) \lesssim c^{-4/3}$, so $\lam{k} = 1 - S(x_k) + O(c^{-2/3})
\approx 1 - c^{-4/3} \approx 1$
--- but in the sense that $1-\lam{k}$ is \emph{not} small
relative to itself, rather $(1-\lam{k}) \sim S(x_k)$).
Using instead the direct estimate:
at $x_k \geq (\log c)^{2/3}$ one has
$\Ai(x_k)^2 \lesssim c^{-4/3}$, and by Proposition~U$'$:
\[
  |\lam{k+d}-\lam{k}| \leq d \cdot c^{-1/3} \cdot \Ai(x_k)^2
  \lesssim c^{1/6} \cdot c^{-1/3} \cdot c^{-4/3} = c^{-3/2}.
\]
Therefore
\[
  P_{k,k+d} \leq c \cdot \frac{c^{-3/2}}{S(x_k)^2}
  \leq c \cdot c^{-3/2} \cdot c^{8/3} \cdot (\log c)^2
  = c^{13/6}(\log c)^2.
\]

\begin{remark}[Direct route for Regime~II]
\label{rem:regimeII_direct}
The above estimate overshoots.
In Regime~II the correct route is not via the denominator
$(1-\lam{k})^{-1}$, but directly:
since $\lam{k} \approx 0$ (not $\approx 1$) for $x_k \gg 0$,
one has $(1-\lam{k})(1-\lam{k+d}) \approx 1$, giving
\[
  P_{k,k+d} \approx c \cdot |\lam{k+d} - \lam{k}|
  \leq c \cdot d \cdot c^{-1/3} \cdot \Ai(x_k)^2
  \leq c^{2/3+1/6} \cdot c^{-4/3}
  = c^{-1/2} \ll c^{1/2}.
\]
\textbf{Regime~II is strongly controlled.} \checkmark
\end{remark}

\subsection{Regime~III: Transition regime,
  $x_k \in [A_0,\, (\log c)^{2/3}]$}
\label{sec:regime3}

This is the critical regime.
In this window:
\begin{itemize}[nosep]
  \item $S(x)$ is positive but decreasing;
  \item $\Phi(x) = \Ai(x)^2/S(x)^2$ is positive and increasing;
  \item the window itself grows with $c$ (its right endpoint
        $(\log c)^{2/3} \to \infty$);
  \item at the right endpoint,
        $\Phi(x_*) \sim x_* \cdot c^{4/3} = (\log c)^{2/3} \cdot c^{4/3}$,
        which grows with $c$.
\end{itemize}
Therefore $\Phi$ is \emph{not} uniformly bounded on this
regime, and the bound
$d \cdot \Phi(\xi_{k,d}) \leq Cc^{-1/6}$ cannot be obtained
from a pointwise upper bound on $\Phi$ alone.

%% ============================================================
\section{Failure-Mode Characterisation}
%% ============================================================

\begin{proposition}[Failure of pointwise methods in Regime~III]
\label{prop:failure}
For $d = 1$ and $x_k \in [A_0, (\log c)^{2/3}]$:
\[
  P_{k,k+1} = c^{2/3}\,\Phi(x_k) + O(c^{1/3}),
\]
and at the right endpoint $x_k \sim (\log c)^{2/3}$:
\[
  c^{2/3}\,\Phi(x_k) \sim c^{2/3} \cdot (\log c)^{2/3} \cdot c^{4/3}
  = c^2 \cdot (\log c)^{2/3} \gg c^{1/2}.
\]
Hence the estimate $P_{k,k+1} \leq Cc^{1/2}$ cannot be derived
from a pointwise upper bound on $\Phi(x_k)$.
\end{proposition}

\begin{remark}[Source of the failure]
\label{rem:failure_source}
The blow-up in Proposition~\ref{prop:failure} is not a
numerical accident.
It reflects the fact that both $\Ai(x)^2$ and $S(x)$ are
exponentially small at $x_k \sim (\log c)^{2/3}$
(both of order $c^{-4/3}$), but their ratio $\Phi(x_k)$
grows polynomially-times-exponentially because $S(x) \sim
\Ai(x)^2/(2x)$ gives $\Phi(x) \sim 4x^2/\Ai(x)^{-2}$
after careful expansion.

Any proof strategy that treats $\Ai(x)^2$ (numerator of $P_{kl}$
after discretisation) and $S(x)^2$ (denominator, from
$(1-\lam{k})^2$) as \emph{independent} quantities discards
the cancellation encoded in the identity $S'(x) = -\Ai(x)^2$,
i.e.\ the fact that numerator and denominator are governed
by the same Airy kernel.
This is Failure-Mode~(b) of Paper~VI (Remark~5.2):
correlated Airy factors.
\end{remark}

\begin{remark}[What is not a failure]
The failure is specific to \emph{pointwise} control of
$\Phi(x_k)$.
It does not preclude B-strong via a \emph{correlated} treatment
of $\Ai(x_k)^2$ and $S(x_k)^2$ as joint functions of $k$.
In particular, the identity
\[
  P_{k,k+d}
  = c^{2/3}\int_{x_k}^{x_{k+d}} \Phi(x)\,dx + \mathcal{E}_{k,d}
\]
(from Remark~\ref{rem:structural}) replaces the mean-value-theorem
bound by an integral, which in principle admits cancellation
between numerator and denominator via the Wronskian structure
of the Airy equation.
Whether this cancellation is sufficient to close B-strong
is the open question identified by this note.
\end{remark}

%% ============================================================
\section{Structural Consequence and Open Problem}
%% ============================================================

\begin{proposition}[Structural consequence]
\label{prop:structural_consequence}
B-strong is equivalent, up to subdominant error terms, to:
\[
  \sup_{\substack{k \leq N,\; d \leq c^{1/6}}}
  \int_{x_k}^{x_{k+d}} \Phi(x)\,dx
  \;\leq\; C\,c^{-1/6}.
\]
This is a \emph{correlated} Airy integral condition on the
Galerkin discretisation scale $c^{-1/3}$:
neither a pointwise bound on $\Phi$,
nor a bound on $\Ai(x)^2$ or $S(x)$ separately,
suffices.
The open problem is precisely whether the Wronskian
identity $S'(x) = -\Ai(x)^2$ can be leveraged to
control the integral $\int_{x_k}^{x_{k+d}} \Ai(x)^2/S(x)^2\,dx$
uniformly in the transition regime.
\end{proposition}

\begin{openproblem}[Closing B-strong via correlated Airy Green function]
\label{op:correlated_airy}
Establish the estimate of Proposition~\ref{prop:structural_consequence}
in the transition regime $x \in [A_0, (\log c)^{2/3}]$.

The natural approach is to replace the pointwise ratio $\Ai(x)^2/S(x)^2$
by an integral kernel that treats $\Ai$ and $S$ as jointly determined
by the Sturm--Liouville equation $y'' = xy$.
Concretely: use the identity
\[
  \int_a^b \frac{\Ai(x)^2}{S(x)^2}\,dx
  = \frac{1}{S(b)} - \frac{1}{S(a)}
\]
and control the right-hand side directly via the WKB
asymptotics of $S(x)$ (not via $\Ai(x)^2$ and $S(x)$
separately).
This requires: explicit WKB expansion of $S(x)$ with
controlled remainder, uniform in $x \in [A_0, (\log c)^{2/3}]$;
and verification that $1/S(x_{k+d}) - 1/S(x_k) \leq Cc^{-1/6}$
for all $k, d$ in the Galerkin regime.

The last inequality is equivalent to B-strong and is the
minimal open statement remaining after this reduction.
\end{openproblem}

%% ============================================================
\section{Summary: Logical Map of the Reduction}
%% ============================================================

\begin{center}
\renewcommand{\arraystretch}{1.4}
\begin{tabular}{llll}
\hline
\textbf{Step} & \textbf{Content} & \textbf{Method} & \textbf{Status} \\
\hline
Telescope & $P_{k,k+d} = c^{2/3}\Delta(1/S) + \mathcal{E}$ & Prop.~U/U$'$ & ✅ exact \\
MVT form & $P_{k,k+d} = c^{2/3}d\,\Phi(\xi) + \mathcal{E}$ & MVT & ✅ valid \\
Regime~I & $|x_k| \leq A_0$: $\Phi = O(1)$, $\delta \leq c^{-1/6}$ & pointwise & ✅ closed \\
Regime~II & $x_k \geq (\log c)^{2/3}$: direct via Prop.~U$'$ & direct & ✅ closed \\
Regime~III & $x_k \in [A_0, (\log c)^{2/3}]$: $\Phi$ not uniform & failure mode & \textbf{⚠ open} \\
Correlated & control $\int \Phi$ via $1/S(b)-1/S(a)$ jointly & Wronskian & \textbf{🔴 open} \\
\hline
\end{tabular}
\end{center}

\medskip\noindent
\textbf{Core diagnosis.}
B-strong does not reduce to separate Airy-amplitude and
eigenvalue-gap estimates.
It requires control of a correlated Airy integral on the
discretisation scale $c^{-1/3}$, where numerator and
denominator are governed by the same Airy kernel via
$S'(x) = -\Ai(x)^2$.
The open problem is whether this Wronskian identity can be
leveraged to prove $1/S(x_{k+d}) - 1/S(x_k) \leq Cc^{-1/6}$
uniformly in the transition regime.

%% ============================================================
\section*{Relation to the Programme}
%% ============================================================

This note is a companion to Paper~VI (\texttt{paper6.tex}).
The two failure modes of Assumption~B-strong identified
in Paper~VI, Remark~5.2 now have a precise quantitative form:
\begin{enumerate}[nosep]
  \item \textbf{Non-uniform stationary phase} (Failure-Mode~(a)):
        corresponds to the non-uniform growth of $\Phi(x)$
        across the transition window.
  \item \textbf{Correlated Airy factors} (Failure-Mode~(b)):
        corresponds to the impossibility of treating
        $\Ai(x)^2$ and $S(x)^2$ independently
        (Remark~\ref{rem:failure_source}).
\end{enumerate}
The reduction of B-strong to the single inequality
$1/S(x_{k+d}) - 1/S(x_k) \leq Cc^{-1/6}$ (Open
Problem~\ref{op:correlated_airy}) is the main new contribution
of this note.
This formulation is the minimal open statement on B-strong
after accounting for all regime-specific controls.

%% ============================================================
\begin{thebibliography}{9}

\bibitem{paper6}
  \textit{Paper~VI: Galerkin Norm Estimates for the Prolate Phase Operator},
  \texttt{paper6.tex}, prolate-gram-coercivity programme, May~2026.

\bibitem{airy_lemma}
  \textit{Airy Discrete Stability Lemma},
  \texttt{airy\_discrete\_stability\_lemma.tex},
  prolate-gram-coercivity programme, 2026.

\bibitem{OsipovRokhlinXiao2013}
  A.~Osipov, V.~Rokhlin, H.~Xiao,
  \textit{Prolate Spheroidal Wave Functions of Order Zero},
  Springer, 2013.

\end{thebibliography}

\end{document}
